\documentclass{article}
\setlength{\emergencystretch}{2em}
\title{A Synthetic Example of Evidence-Bounded Result Writing}
\author{Super Writer Tutorial}
\date{}
\begin{document}
\maketitle
\section{Scope}
This document uses hand-authored synthetic teaching data. It is not a research
finding and makes no claim of statistical significance or benchmark leadership.
\section{Protocol}
Each method is assigned 100 synthetic trials under clean and perturbed input.
Success is reported as a count divided by total trials. Per-trial outcomes,
uncertainty estimates, and runtime measurements are unavailable.
\section{Results}
Table~\ref{tab:synthetic} records all observations in the example. The candidate
has descriptive success rates of 84\% and 69\%, compared with 82\% and 61\%
for the baseline. These differences are 2 and 8 percentage points.
\begin{table}[ht]
\centering
\begin{tabular}{lrr}
Condition & Baseline & Candidate \\
Clean & 82/100 & 84/100 \\
Perturbed & 61/100 & 69/100 \\
\end{tabular}
\caption{Synthetic teaching counts, not real experimental results.}
\label{tab:synthetic}
\end{table}
\section{Limitations}
The example supports only a descriptive comparison of two synthetic conditions.
It does not establish significance or causal effects. Performance in other
conditions and possible costs remain unknown.
\end{document}
